\section{Three-state recurrent-cycle realization}\label{sec:realization}

A recurrent word must satisfy both capacity and shared-table compatibility.
We characterize the latter on three-state paths; initial-state accessibility remains
a separate question. Figure~\ref{fig:realization-pipeline} previews the construction.

\subsection{The normal form and its scope}

\begin{proposition}[Three-state path normal form]\label{prop:three-normal-form}
Let a simple three-state tagged cycle have used path support with $n\ge3$ vertices.
Every edge multiplicity is one or two; the doubled edges form a matching.
Up to cyclic shift its spatial word is the ordinary contour, with one passage of
each doubled edge replaced by the three-step passage ``forward, back, forward.''
The choice of the forward or backward contour passage is its insertion phase.
\end{proposition}
\begin{proof}
For adjacent edges, capacity gives $k_{j-1}+k_j\le3$ and both multiplicities are
positive. Hence every $k_j$ is one or two, and two doubled edges cannot be adjacent.
A neighbouring single edge is crossed once in each direction. The additional two
crossings of a doubled edge must therefore lie within one of its two contour
passages; otherwise a neighbouring single-edge cut would be crossed again.
They form the stated immediate backtrack. Endpoint insertions can have equivalent
cyclic descriptions, so uniqueness of the parametrization is not asserted.
\end{proof}

For a one-edge support the multiplicity can instead be three; that case is not a
matching input and is treated separately in Lemma~\ref{lem:one-edge-realization}.
The normal-form necessity concerns used support and remains valid in a larger maze.
The converse below uses the masks of the path itself. A proper interval of an ambient
path requires the boundary-aware conditions in Appendix~\ref{app:boundary-aware}.

Fix a finite induced grid path $P=(v_0,\ldots,v_{n-1})$, $n\ge2$, with its absolute
compass orientation and exactly its own masks. Choose a matching
$\mathcal D_2\subseteq\{0,\ldots,n-2\}$ and phases $\eta_e\in\{0,1\}$.
In the contour $0,1,\ldots,n-1,n-2,\ldots,1$, replace the index-increasing passage
of $e$ if $\eta_e=0$, and its index-decreasing passage if $\eta_e=1$, by the
three-step backtrack. Denote this specification by $S$, and write
\[
 I=\Z/L\Z,\quad L=2(n-1)+2|\mathcal D_2|,\quad
 w_i\in V(P),\quad \sigma(i)=i+1,\quad M_i=M_P(w_i).
\]
Let $d_i$ be the direction from $w_i$ to $w_{\sigma(i)}$.
An occurrence means an index $i\in I$, not just its cell. Internal cells have load
two or three, and endpoints load one or two. An internal cell of load three is
\emph{saturated}. Realization asks for a legal three-state table and states $q_i$
such that the $L$ distinct tags $(w_i,q_i)$ follow this cyclic word. No primitive
spatial-word hypothesis or state-zero accessibility requirement is imposed.

A necessary local condition is \emph{same-mask doubled-side consistency}: all
saturated cells of one absolute mask must use the same port twice. We call this the
direction condition. The finite profile family and matrices in the next theorem are
constructed explicitly below from masks, occurrences and successor links.

\begin{theorem}[Exact three-state realization]\label{thm:three-state-realization}
For every specification $S$ just defined, the following are equivalent:
\begin{enumerate}[label=(\roman*),leftmargin=2em]
\item $S$ is realized by a simple three-state tagged cycle on the path with its own masks;
\item the direction condition holds and some $\theta\in\Theta(S)$ satisfies
\[
 \rank_{\Ftwo}A_\theta=\rank_{\Ftwo}[A_\theta\mid b_\theta];
\]
\item the direction condition holds and there is \emph{one} $\theta\in\Theta(S)$ such that
\[
 \forall\lambda:\quad\lambda^TA_\theta=0\ \Longrightarrow\ \lambda^Tb_\theta=0.
\]
\end{enumerate}
From a consistent profile and its solution one constructs all $45$ legal controller
rows. The resulting cycle has least tagged period exactly
$L=2(n-1)+2|\mathcal D_2|$. Every realizing state assignment belongs to at least one
profile. The same profile is used throughout all components of the specification.
\end{theorem}

% A diagram of the existing finite-profile construction. All components share theta.
\begin{figure}[tbp]
\centering
\begin{tikzpicture}[maze,x=1mm,y=1mm]
 \foreach \x in {0,42,84,126}{\draw[lightpanel] (\x,17) rectangle (\x+34,42);}
 \node[font=\footnotesize\sffamily\bfseries,align=center] at (17,36) {SPECIFICATION};
 \draw[edge] (6,25)--(13,25)--(13,29)--(23,29)--(29,29);
 \draw[supportA] (6,25)--(13,25); \draw[supportA] (23,29)--(29,29);
 \foreach \x/\y in {6/25,13/25,13/29,23/29,29/29}{\node[cell] at (\x,\y) {};}
 \node[font=\scriptsize] at (17,20) {path, matching, phases};
 \node[font=\footnotesize\sffamily\bfseries,align=center] at (59,36) {ONE PROFILE $\theta$};
 \node[font=\footnotesize,align=center] at (59,25) {roles; endpoint states\\differences $\kappa_M$};
 \node[font=\footnotesize\sffamily\bfseries,align=center] at (101,36) {AFFINE SYSTEM};
 \node[font=\small] at (101,27) {$A_\theta z=b_\theta$};
 \node[font=\scriptsize] at (101,20) {one bit per internal cell};
 \node[font=\footnotesize\sffamily\bfseries,align=center] at (143,36) {CONTROLLER};
 \node[font=\footnotesize,align=center] at (143,25) {consistent used rows\\+ legal completion};
 \foreach \x in {35,77,119}{\draw[flow] (\x,29.5)--(\x+6,29.5);}
 \node[font=\footnotesize,anchor=west] at (1,8) {Repeated mask $M=\N\W$, paired port $\W$:};
 \node[font=\footnotesize,anchor=west] at (90,8) {$Y_i+Y_j=\kappa_M(t_i+t_j)$};
 \node[font=\footnotesize,anchor=west] at (1,0) {$t_i=t_j$ forces the same successor.};
 \node[font=\footnotesize,anchor=west] at (83,0) {$\kappa_M=0$ allows merged outputs.};
\end{tikzpicture}
\caption{The realization construction, with $Y_i=X_{\sigma(i)}$. The same profile serves every repeated-mask component; varying the profile gives a finite union of affine solution sets, not one XOR system.}
\label{fig:realization-pipeline}
\end{figure}


\subsection{Two-input functionality and the finite profile}

The algebraic mechanism concerns a function on two inputs, not an invertible
successor map. Encode the three states by the nonzero vectors
\[
 H=\Ftwo^2,\qquad Q_3=H\setminus\{0\}=\{1,2,3\},
\]
where the numbers are binary vector codes and addition is XOR. The final controller
names are obtained by $1\mapsto0$, $2\mapsto1$, $3\mapsto2$.

\begin{lemma}[Two-input successor difference]\label{lem:successor-difference}
For input bits $t_i\in\Ftwo$ and output states $Y_i\in Q_3$, the conditions
\[
 t_i=t_j\ \Longrightarrow\ Y_i=Y_j\quad\text{for all }i,j
\]
are equivalent to the existence of $\kappa\in H$ with
\[
 Y_i+Y_j=\kappa(t_i+t_j)\quad\text{for all }i,j.
\]
\end{lemma}
\begin{proof}
If only one input occurs, all outputs coincide and $\kappa=0$ works. If both occur,
functionality gives two outputs $Y_0,Y_1$; take $\kappa=Y_0+Y_1$.
Conversely, equal inputs make the right side zero and force equal outputs.
The case $Y_0=Y_1$ is allowed and is exactly the zero difference.
\end{proof}

For two occurrences of $M=\N\W$ on its paired $\W$-side, take $\kappa=3=(1,1)$ and role bits
$t_i=z_u$, $t_j=z_v$. If the observed successors are $Y_i=1$, $Y_j=2$,
the lemma gives $3=3(z_u+z_v)$, hence $z_u+z_v=1$ over $\Ftwo$.
Equal outputs instead require equal bits for this nonzero difference; $\kappa=0$
allows both inputs to share an output. This is an algebraic illustration, not a
claim that a separately prescribed path is realizable.

Every internal vertex uses both ports. A realizing three-state direction map at its
mask is therefore a surjection with fibres of sizes one and two. For each internal
mask $M$ choose its singleton port $u_M$, its singleton state $a_M\in Q_3$, and let
$p_M$ be the other port. If $M$ has a saturated representative, $p_M$ must be its
doubled port, so $u_M$ is forced; otherwise both port choices are retained.
Choose a successor difference $\kappa_M\in H$, restricted to $\kappa_M\ne0$ when
$M$ is saturated. For each endpoint separately choose an injective assignment of
its actual occurrences to $Q_3$. These choices form $\theta\in\Theta(S)$.
They assign no arbitrary row outputs.

Let $b(a)$ be the smaller numeric code in $Q_3\setminus\{a\}$.
For every internal cell $v$ introduce one bit $z_v$. Define its occurrence states by
\begin{equation}\label{eq:real-roles}
 X_i=\begin{cases}
 a_{M_i},&d_i=u_{M_i},\\
 b(a_{M_i})+a_{M_i}t_i,&d_i=p_{M_i},
 \end{cases}
 \qquad t_i=z_{w_i}+\epsilon_i.
\end{equation}
At a load-two cell the unique paired-side occurrence has $\epsilon_i=0$.
At a saturated cell the two such occurrences have $\epsilon=0,1$ in temporal
order after choosing an origin of the cyclic word. Endpoint $X_i$ are the constants
chosen in the profile. The two paired states are $b(a_M)$ and $b(a_M)+a_M$,
both nonzero and different from $a_M$. Thus local injectivity is built in.
The state absent over a load-two cell is still a possible row state of its mask;
it cannot be discarded globally.

\subsection{The canonical affine system}

For an internal mask form
\[
 U_M=\{i:M_i=M,\ d_i=u_M\},\qquad
 B_M=\{i:M_i=M,\ d_i=p_M\}.
\]
Both groups are nonempty. In each choose an anchor $j$, for example the least
occurrence index. For every other occurrence $i$ impose respectively
\begin{align}
 X_{\sigma(i)}+X_{\sigma(j)}&=0 &&(i,j\in U_M),\label{eq:real-U}\\
 X_{\sigma(i)}+X_{\sigma(j)}&=\kappa_M(t_i+t_j)
 &&(i,j\in B_M).\label{eq:real-B}
\end{align}
Endpoint occurrences are grouped by the \emph{pair} consisting of their one-port
mask and their chosen state. For every non-anchor member of such a group impose
\begin{equation}\label{eq:real-E}
 X_{\sigma(i)}+X_{\sigma(j)}=0.
\end{equation}
In particular, equal endpoint masks cannot be ignored. Equal states at different
masks do not identify rows, and different states at the same endpoint mask are not
identified either.

For a fixed $\theta$, all $a_M$, $\kappa_M$ and endpoint assignments are constants.
Expanding the vector equations into their two scalar coordinates gives
\begin{equation}\label{eq:real-system}
 A_\theta z=b_\theta,\qquad z\in\Ftwo^{n-2}.
\end{equation}
Constant equations are retained: $0=1$ rejects a profile. For $n=2$ there are no
internal bits and the same construction is a zero-dimensional system.

A group of $h$ occurrences contributes $h-1$ vector comparisons. There are fewer
than $2L$ scalar rows, each involving at most two source bits and two successor bits.
The symbolic object consequently has size $O(n)$ and is constructed without solving
realizability. If $m_3$ internal masks have saturated representatives and $m_2$ do
not, then $m_3+m_2\le6$ and
\begin{equation}\label{eq:profile-bound}
 |\Theta(S)|\le36\,9^{m_3}24^{m_2}\le36\,24^6.
\end{equation}
This deliberately coarse constant is independent of the path length. Profiles may
overlap; uniqueness of representation is not required.

\subsection{Proof of the realization theorem}

\begin{proof}[Proof of Theorem~\ref{thm:three-state-realization}]
\emph{Necessity.}
Let a legal table realize the distinct tags. Encode its states in $Q_3$.
At every internal mask both directions occur, so its direction fibres have sizes
one and two. At a saturated cell all three states occur, with the two states on the
doubled side distinct. Therefore every saturated representative of that mask doubles
the same side. Extract $u_M,a_M,p_M$ from the table. At each cell the paired state
or ordered pair determines a bit $z_v$ in \eqref{eq:real-roles}. The endpoint states
give the required injections.

All sources in $U_M$ have the same state and mask; deterministic successor states
give \eqref{eq:real-U}. In $B_M$, source-state equality is equivalent to equality
of $t_i$. Lemma~\ref{lem:successor-difference}, with $Y_i=X_{\sigma(i)}$,
gives a $\kappa_M$ satisfying \eqref{eq:real-B}.
When the mask is saturated, its two paired occurrences at one cell move to the
same neighbour along the doubled edge. Their successor occurrences are distinct,
and distinct tags over that neighbour have different states. Hence their difference
is nonzero, as required in the profile. This also covers a doubled edge incident to
an endpoint. Equal endpoint keys give exactly \eqref{eq:real-E}.
Thus the extracted profile and bits solve \eqref{eq:real-system}; every row
dependency has zero right-hand-side sum.

These comparisons exhaust repeated keys. Different masks mean different rows.
For the same internal mask, singleton and paired states are disjoint, and each
fibre is handled above. A one-port endpoint mask cannot equal a two-port internal
mask. The remaining endpoint coincidences are precisely the groups in
\eqref{eq:real-E}. In particular no additional successor-row condition is missing.

\emph{Linear consistency.}
Row elimination over $\Ftwo$ produces an inconsistent equation exactly when a sum
of original rows has zero left side and right side one. If no such row occurs,
assign the free bits and back-substitute the pivot bits. This proves the equivalence
of rank consistency, the stated balance condition, and existence of a solution,
including the zero-dimensional case.

\emph{Sufficiency and reconstruction.}
Take one consistent profile and a solution. Equations~\eqref{eq:real-roles}
produce nonzero state vectors. At a load-two cell the singleton and paired states
are different. At a saturated cell the opposite $\epsilon$ values produce both
paired states, distinct from each other and from the singleton. Endpoint injections
are part of the profile. Hence all $L$ resulting tags are distinct.

Apply the one global renaming to actual controller states $q_i\in\{0,1,2\}$ and
prescribe every used row by
\begin{equation}\label{eq:real-used-controller}
 \delta(q_i,M_i)=(d_i,q_{\sigma(i)}).
\end{equation}
Suppose two prescriptions have the same key. At an internal mask both sources lie
in the same direction fibre. In the singleton fibre, adding their anchor equations
gives equal successors. In the paired fibre, equal states give $t_i=t_j$, and
adding their anchor equations yields
\[
 X_{\sigma(i)}+X_{\sigma(j)}=\kappa_M(t_i+t_j)=0.
\]
The direction already agrees. At an endpoint the direction is unique, and
\eqref{eq:real-E} gives equal successors for equal keys. Thus
\eqref{eq:real-used-controller} is a function, not a conflicting set of assignments.

Every used direction belongs to its mask. For each still unassigned key $(q,M)$,
$M\ne\varnothing$, choose any $d\in M$ and any next state. This produces all
$3(2^4-1)=45$ legal rows without changing a used row. If only one paired input was
observed, its profile difference need not prescribe the absent row's output.
In particular, one must not extrapolate it to a zero vector and treat that vector
as a state. The observed difference may instead be replaced by zero without
changing the realized assignments.

Starting at $(w_0,q_0)$, the prescribed row takes the controller to
$(w_{\sigma(0)},q_{\sigma(0)})$; induction follows the entire specified word.
In each selected compass direction the next cell is the prescribed neighbour.
The orbit returns after $L$ steps, and the preceding $L$ tags are distinct.
Its least tagged period is therefore exactly $L$. This argument neither assumes
nor needs that the spatial projection have least period $L$.
\end{proof}

The same key comparisons close all successor equalities. Within a fixed profile,
linear elimination tests or projects partial assignments; across profiles, one takes
a union. Unsaturated vertices therefore have no automatic completion rule.

\subsection{Why one XOR system is not enough}

Consider the fixed specification
\[
 \text{directions }\E\N\E\N\N\N\E,\qquad
 \mathcal D_2=\{0,2,6\},\qquad \eta_e=0,
\]
whose masks are $\E,\N\W,\E\Sdir,\N\W,\N\Sdir,\N\Sdir,\E\Sdir,\W$.
Its cyclic word is
\[
 (0,1,0,1,2,3,2,3,4,5,6,7,6,7,6,5,4,3,2,1).
\]
Let $x$ denote the same/swap correspondence of the ordered west-state pairs
$(q_1,q_{19})$ and $(q_5,q_{17})$ at the two saturated $\N\W$ cells.
Let $y$ denote that of the east-state pairs $(q_4,q_6)$ and $(q_{10},q_{12})$
at the two saturated $\E\Sdir$ cells; zero means same and one swap.
The two intervening $\N\Sdir$ cells are unsaturated.

\begin{proposition}[A non-affine saturated correspondence relation]\label{prop:nand8}
For this fixed spatial specification, the exact extendible relation is
\[
 \operatorname{Ext}(x,y)=\{00,01,10\}.
\]
\end{proposition}
\begin{proof}
The singleton $\N\W/\N$ rows give $q_3=q_7$ and $q_4=q_8$.
The singleton $\E\Sdir/\Sdir$ rows give $q_{18}=q_{14}$ and $q_{19}=q_{15}$.
If $x=1$, then $q_{17}=q_1\ne q_{19}$. Equality of the two $\N\Sdir/\Sdir$
source states, $q_{16}=q_{15}$, would force $q_{17}=q_{16}=q_{19}$ by their
successor rows. Thus two different south states are needed. If $y=1$, then
$q_{10}=q_6\ne q_4=q_8$. Equality of the north states $q_8=q_9$ would force
$q_9=q_{10}$, again a contradiction. Hence two different north states are needed.
North and south states at the same mask cannot coincide. The assignment $11$
would therefore need four states. Appendix~\ref{app:nand-witnesses} gives full
realizing state assignments for the other three choices.
\end{proof}

An XOR solution set and its projections are affine, with one, two or four possible
assignments on two bits, never three. Even auxiliary XOR variables therefore cannot
describe these actual correspondences by one system. This excludes faithful
correspondence encodings, not arbitrary hypothetical decision-only graphs. The
insertion phases $\eta_e$ are fixed, distinct from $x,y$; the unpinned specification
is realizable.

\subsection{An infinite phase-locking family}\label{subsec:phase-locking}

\begin{theorem}[Phase locking across arbitrary gaps]\label{thm:phase-locking}
Let $r\ge1$ and $\ell_0,\ldots,\ell_r\ge1$. On the induced monotone path
\[
 \E^{\ell_0}\N\E^{\ell_1}\N\cdots\N\E^{\ell_r},
\]
double all $r$ north edges and no others, assigning phases $\eta_1,\ldots,\eta_r$.
The specification is realizable with three states if and only if all phases are equal.
\end{theorem}
\begin{proof}
The lower endpoint of each doubled north edge has mask $\N\W$ and singleton west
state $w$; the upper endpoint has mask $\E\Sdir$ and singleton east state $e$.
At the $j$th edge denote its temporal north states by $A_j,B_j$ and its south
states by $C_j,D_j$. Then
$\{A_j,B_j,w\}=Q$ and $\{C_j,D_j,e\}=Q$.
For either insertion phase the south rows give
\[
 C_j\longmapsto B_j,\qquad D_j\longmapsto w.
\]
Two pairs $(C_j,D_j)$ and $(C_k,D_k)$ either agree or are swapped.
A swap would identify $B_j$ with $w$, so the pairs agree; then $B_j=B_k$ and
$A_j=A_k$. The north rows are
\[
\begin{array}{c|cc}
\eta_j& A_j & B_j\\\hline
0&C_j&e\\
1&e&D_j
\end{array}
\]
and distinct phases would identify a paired state with the singleton $e$.
Thus all phases are equal, independently of the horizontal gaps.

For sufficiency, use the following rows; the two rows for each phase list the
three outputs in state order $0,1,2$:
\[
\begin{array}{c|ccc}
 (\eta,M)&0&1&2\\\hline
 (0,\N\W)&(\W,2)&(\N,1)&(\N,0)\\
 (0,\E\Sdir)&(\E,1)&(\Sdir,2)&(\Sdir,0)\\
 (1,\N\W)&(\W,2)&(\N,0)&(\N,1)\\
 (1,\E\Sdir)&(\E,1)&(\Sdir,0)&(\Sdir,2)
\end{array}
\]
For either phase add
$\delta(1,\E\W)=(\E,1)$, $\delta(2,\E\W)=(\W,2)$,
$\delta(2,\E)=(\E,1)$, and $\delta(1,\W)=(\W,2)$.
Start at the left endpoint in state $2$. At phase zero the forward doubled-edge
sequence is $\N\W_1,\E\Sdir_1,\N\W_2,\E\Sdir_0$, with backward sequence
$\E\Sdir_2,\N\W_0$. At phase one the forward sequence is
$\N\W_1,\E\Sdir_0$, and the backward sequence is
$\E\Sdir_2,\N\W_2,\E\Sdir_1,\N\W_0$.
Here a mask with a subscript denotes the occurrence at the corresponding endpoint
of that edge with that state. Horizontal segments preserve state $1$ eastward and
state $2$ westward. The endpoint rows close the word; all occurrences at one cell
have different states. Completing unused rows gives the required simple cycle.
\end{proof}

With $r=2$ and all gaps one, unequal phases give the six-vertex staircase.
Its contradiction is equivalently the three constraints $x=1$, $x+y=0$, $y=0$
from the general successor comparisons, whose sum is $0=1$; an explicit derivation
is recorded in Appendix~\ref{app:staircase}. Minimality of six among residual
obstructions is a finite-check statement, not needed for this proof.
More generally, for any fixed path radius $R$, choose gaps $(1,h,1)$ with
$h>2R+2$ and phases $(0,1)$. Every radius-$R$ neighbourhood occurs in a realizable
uniform-phase specification, but the whole specification is unrealizable.
Thus bounded geometric neighbourhoods do not replace the global mask constraints.
This is not a lower bound on the size of an algebraic certificate.

\subsection{Recognition, reconstruction and the accessibility boundary}

Enumerate the finite profiles after checking the direction condition, compile
\eqref{eq:real-system}, and solve by Gaussian elimination. A successful branch
constructs the table by \eqref{eq:real-used-controller}; rejection requires failure
of every admissible profile. With $K$ bounded as in \eqref{eq:profile-bound}, this
gives a coarse $O(Kn^3)$ time bound and $O(n^2)$ working space. Fixed-alphabet
polynomiality is not asserted as a separate novelty: finitely many full tables already
exist independently of $n$. The content is the special sparse normal form and
constructive removal of arbitrary row outputs.

Realization does not determine state-zero accessibility: two completions on one
four-vertex path preserve the same contour while making it accessible from every
state-zero start or from none. Appendix~\ref{app:basins} gives both tables and a
conditional uniqueness criterion. Universal three-state rendezvous remains open.
